% Generated draft: requires mathematical and source review.
\documentclass[11pt]{article}
\usepackage[margin=25mm]{geometry}
\usepackage{fontspec}
\setmainfont{Latin Modern Roman}
\usepackage{amsmath,amssymb,mathtools,unicode-math}
\setmathfont{Latin Modern Math}
\usepackage{enumitem,needspace,xurl,hyperref}
\setlength{\parindent}{0pt}
\setlength{\parskip}{5pt}
\setlength{\emergencystretch}{3em}
\newcommand{\N}{\mathbb{N}}
\newcommand{\Z}{\mathbb{Z}}
\newcommand{\Q}{\mathbb{Q}}
\newcommand{\R}{\mathbb{R}}
\newcommand{\C}{\mathbb{C}}
\newcommand{\F}{\mathbb{F}}
\newcommand{\A}{\mathbb{A}}
\newcommand{\PP}{\mathbb{P}}
\newcommand{\E}{\mathbb{E}}
\newcommand{\eps}{\varepsilon}
\newcommand{\dd}{\,\mathrm{d}}
\begin{document}
\section*{UnsolvedMath ID: 30004132}
\textbf{Source problem code:} catalogue entry
\subsection{Definitions and mathematical statement}
For a finite real symmetric edge-weight matrix \(B\) of size \(q\geq 1\), write its eigenvalues as \(\lambda_1(B)\geq\cdots\geq\lambda_q(B)\). The supplied record calls a weighted graph \(H\) \emph{antiferromagnetic} when its edge-weight matrix has at most one strictly positive eigenvalue; in this matrix setting, this means
\[
\#\{i:\lambda_i(B)>0\}\leq 1.
\]
Equivalently, \(\lambda_i(B)\leq 0\) for every \(2\leq i\leq q\); the condition does not require \(\lambda_1(B)>0\). A graph is \(d\)-regular if every vertex has degree \(d\), and triangle-free if it contains no cycle of length three. The notation \(\operatorname{hom}(G,H)\) denotes the weighted homomorphism quantity mentioned in the source, whose precise weight convention is not supplied.

The proposed extension asks whether, for every admissible antiferromagnetic target \(H\) and every admissible degree \(d\), the reverse-Sidorenko homomorphism inequality referred to in the record holds for every \(d\)-regular source graph \(G\), with no triangle-free assumption.

This is an unresolved formulation of the intended question: the record does not state the inequality, its comparison quantity, its normalization, or its complete hypotheses. It also does not specify the allowed edge and vertex weights, loop conventions, finiteness assumptions, or degree range. The finite real symmetric matrix setting above makes the supplied spectral condition precise, but is not independently confirmed as the full intended target class. A complete quantified inequality cannot be recovered from this record alone.
\subsection{Short English statement}
The question asks whether the stated antiferromagnetic spectral condition on a weighted target \(H\) permits removing triangle-freeness from a reverse-Sidorenko bound for \(d\)-regular source graphs. The actual bound and several graph and weight conventions are absent, so its precise mathematical formulation requires source review. The supplied record labels the problem open.
\subsection{Sources}
\begin{itemize}
\item Local registry: \texttt{data/unsolvedmath/problems.json}, with \texttt{catalog/problems.json} as the additional registry specified by the project; permanent UnsolvedMath ID \(30004132\), source problem code \texttt{catalogue entry}. This draft uses the supplied complete source record; no registry files or external sources were independently inspected.
\item Supplied report citation: \emph{Combinatorics, Probability and Computing (2020)}, Oberwolfach Reports, DOI \texttt{10.4171/owr/2019/19}, pp.\ 1113--1189. The same reference is listed among the literature sources as ``Oberwolfach Report 19/2019, weighted homomorphisms'': \url{https://doi.org/10.4171/owr/2019/19}. The record gives proposed year 2019 and names no proposer.
\item Supplied literature source: Lee--Oh--Seo, \emph{Counting homomorphisms in antiferromagnetic graphs via Lorentzian polynomials}: \url{https://arxiv.org/abs/2506.13659}.
\item Literature assessment dated exactly \texttt{2026-08-22}: the supplied curation labels the problem open and states that the antiferromagnetic extension remains conjectural in the original reverse-Sidorenko paper, while 2025 Lorentzian-polynomial work establishes important subclasses but not all antiferromagnetic targets and source graphs. This is supplied dated curation, not an independently verified literature assessment. The record requests independent verification of time-sensitive status and cited claims.
\item The supplied statement status is \texttt{corrected\_verified}, with an explanation that mathematical content was preserved while notation, typography, and sentence structure were repaired. This supplied editorial designation does not establish mathematical correctness or completeness.
\end{itemize}
\end{document}
